\documentclass[11pt]{article}
\usepackage[utf8]{inputenc}
\usepackage[T1]{fontenc}
\usepackage[brazilian]{babel}
\usepackage[a4paper,margin=1.4cm]{geometry}
\usepackage[scaled=.92]{helvet}
\renewcommand{\familydefault}{\sfdefault}
\usepackage{amsmath}
\usepackage{booktabs}
\usepackage[table]{xcolor}
\usepackage{tikz}
\usepackage{tcolorbox}
\tcbuselibrary{skins}

\definecolor{pena}{HTML}{2F6FDE}
\definecolor{folha}{HTML}{2E9E6A}
\definecolor{brasa}{HTML}{E0663A}
\definecolor{tinta}{HTML}{1F2937}
\definecolor{nevoa}{HTML}{6B7280}

\pagestyle{empty}

\begin{document}

\begin{tcolorbox}[colback=pena, colframe=pena, sharp corners]
  \color{white}
  {\LARGE\bfseries Observatório do Rio Claro}

  {\large Boletim de qualidade da água --- 3º trimestre de 2026}

  {\small Laboratório de Hidrologia \quad$\cdot$\quad 18 de outubro de 2026}
\end{tcolorbox}
\medskip

\begin{center}
\begin{tabular}{cccc}
  {\Huge\bfseries\color{pena} 4} & {\Huge\bfseries\color{brasa} $-$18\%} & {\Huge\bfseries\color{folha} 6,3} & {\Huge\bfseries\color{tinta} 41} \\
  {\small\color{nevoa} pontos de coleta} & {\small\color{nevoa} vazão média} & {\small\color{nevoa} menor O$_2$ (mg/L)} & {\small\color{nevoa} maior turbidez (NTU)} \\
\end{tabular}
\end{center}

\section{Resumo}

Este boletim reúne as medições feitas em quatro pontos do \textbf{Rio Claro} entre julho e setembro.
A vazão média caiu \textcolor{brasa}{\textbf{18\%}} em relação ao trimestre anterior, efeito direto
da estiagem, enquanto o oxigênio dissolvido se manteve dentro do limite da \emph{classe~2} em todos
os pontos. A turbidez subiu no trecho urbano, onde a concentração de sedimentos acompanha as obras
de drenagem iniciadas em agosto.

\section{Medições por ponto}

\begin{center}
\rowcolors{2}{white}{pena!6}
\begin{tabular}{lrrr}
  \toprule
  \textbf{Ponto} & \textbf{Vazão (m$^3$/s)} & \textbf{O$_2$ (mg/L)} & \textbf{Turbidez (NTU)} \\
  \midrule
  P1 --- Nascente    & 2,4 & 8,1 & 3  \\
  P2 --- Ponte Velha & 5,7 & 7,4 & 12 \\
  P3 --- Centro      & 6,2 & 6,3 & \textcolor{brasa}{\textbf{41}} \\
  P4 --- Foz         & 7,9 & 6,8 & 27 \\
  \bottomrule
\end{tabular}
\end{center}

\begin{center}
\begin{tikzpicture}[x=1.1cm, y=0.026cm]
  \draw[dashed, gray] (0,25) -- (8.8,25);
  \draw[dashed, gray] (0,50) -- (8.8,50);
  \draw[dashed, gray] (0,75) -- (8.8,75);
  \node at (-0.4,25) {2,5}; \node at (-0.4,50) {5,0}; \node at (-0.4,75) {7,5};
  \draw[fill=pena, draw=pena]  (0.6,0) rectangle (1.8,24); \node at (1.2,31) {\textbf{2,4}};
  \draw[fill=pena, draw=pena]  (2.6,0) rectangle (3.8,57); \node at (3.2,64) {\textbf{5,7}};
  \draw[fill=brasa, draw=brasa] (4.6,0) rectangle (5.8,62); \node at (5.2,69) {\textbf{6,2}};
  \draw[fill=folha, draw=folha] (6.6,0) rectangle (7.8,79); \node at (7.2,86) {\textbf{7,9}};
  \draw[thick] (0,0) -- (8.8,0);
  \node at (1.2,-8) {P1}; \node at (3.2,-8) {P2}; \node at (5.2,-8) {P3}; \node at (7.2,-8) {P4};
\end{tikzpicture}\\[2pt]
{\small\color{nevoa} Vazão média de cada ponto (m$^3$/s).}
\end{center}

\section{Modelo de vazão}

A vazão $Q$ em cada seção vem da fórmula de Manning, com a área molhada $A$, o raio hidráulico $R$,
a declividade $S$ e a rugosidade $n$; a carga de sedimentos $M$ é a integral da vazão pela concentração $C$:
\begin{equation}
  Q = \frac{1}{n}\, A\, R^{2/3} \sqrt{S}
  \qquad\qquad
  M = \int_{t_0}^{t_1} Q(t)\, C(t)\, \mathrm{d}t \approx \sum_{k=1}^{N} Q_k\, C_k\, \Delta t
\end{equation}

\begin{tcolorbox}[colback=brasa!8, colframe=brasa, title=Destaque do trimestre, fonttitle=\bfseries]
  O ponto P3 passou de 25 para \textbf{41 NTU}. A equipe recomenda uma nova coleta em novembro,
  depois do fim das obras, para separar o efeito da estiagem do efeito da drenagem.
\end{tcolorbox}

\section{Próximos passos}

\begin{enumerate}
  \item Instalar o sensor contínuo de turbidez no ponto P3.
  \item Comparar a série de 2026 com a média dos últimos dez anos.
  \item Publicar os dados abertos no portal do município.
\end{enumerate}

\end{document}
